\section{Track H: Computational Methods for Low-discrepancy Sampling and Applications}\newpage

\begin{talk}
  {Searching Permutations for Constructing Low-Discrepancy Point Sets and Investigating the Kritzinger Sequence}% [1] talk title
  {Fran\c{c}ois Cl\'{e}ment}% [2] speaker name
  {Department of Mathematics, University of Washington}% [3] affiliations
  {fclement@uw.edu}% [4] email
  {Carola Doerr, Kathrin Klamroth, Lu\'is Paquete}% [5] coauthors
  
    
   
This talk focuses on two different approaches for the construction of low-discrepancy sets that are quite different from traditional approaches, yet yield excellent empirical results in two dimensions. The first of these, which will be the main focus of my talk, is based on selecting the relative position of the different points we wish to place, before using non-linear programming methods to obtain a point set with extremely low star discrepancy. In [1], we showed that this method consistently outperformed all other existing techniques. It is however subject to computational limits: finding good permutation choices, with or without optimization, is the next key step in improving our understanding of low-discrepancy structures.

In the second part of the talk, I will quickly highlight some extended numerical experiments on the sequence introduced by Kritzinger in [2], showing that despite the lack of theoretical results proving that it is a low-discrepancy sequence, it performs at least as well as known sequences in one dimension, despite being constructed greedily. 

\medskip

\begin{enumerate}
 \item[{[1]}] F. Cl\'ement, C. Doerr, K. Klamroth, L.Paquete (2024). {\it Transforming the Challenge of Constructing Low-Discrepancy Point Sets into a Permutation Selection Problem}., to appear, arxiv: https://arxiv.org/abs/2407.11533.
 \item[{[2]}] R. Kritzinger (2022).  {\it Uniformly Distributed Sequence generated by a greedy minimization of the $L_2$ discrepancy}., Moscow Journal of Combinatorics and Number Theory, \textbf{11}(2), 215--236.
\end{enumerate}

\end{talk}

\begin{talk}
  {Minimizing the Stein Discrepancy}% [1] talk title
  {Nathan Kirk}% [2] speaker name
  {Illinois Institute of Technology}% [3] affiliations
  {nkirk@iit.edu}% [4] email
  {}% [5] coauthors
  {Computational Methods for Low-discrepancy Sampling and Applications}% [6] special session. Leave this field empty for contributed talks. 
    
   
Approximating a probability distribution using a discrete set of points is a fundamental task in modern scientific computation, with applications in uncertainty quantification. Points whose empirical distribution is close to the true distribution are called low-discrepancy. We discuss recent advances in this area, including the use of Stein discrepancies and various optimization techniques. In particular, we introduce Stein-Message-Passing Monte Carlo (Stein-MPMC) [1], an extension of the original Message-Passing Monte Carlo model and the first machine-learning algorithm for generating low-discrepancy (space-filling) point sets. Additionally, we present a generalized Subset Selection algorithm [2], a much simpler yet highly effective optimization method.

\medskip

\begin{enumerate}
 \item[{[1]}] N. Kirk, T. K. Rusch, J. Zech, D. Rus, \textit{Low Stein Discrepancy via Message-Passing Monte Carlo}, Preprint (2025)
 \item[{[2]}] D. Chen, F. Cl\'ement, C. Doerr, N. Kirk, L. Paquette, \textit{Optimizing Kernel Discrepancies via Subset Selection}, Preprint (2025)
\end{enumerate}


\end{talk}

\begin{talk}
  {Improving Efficiency of Sampling-based Motion Planning via Message-Passing Monte Carlo}% [1] talk title
  {Makram Chahine}% [2] speaker name
  {CSAIL, Massachusetts Institute of Technology}% [3] affiliations
  {chahine@mit.edu}% [4] email
  {T. Konstantin Rusch, Zach J. Patterson, and Daniela Rus}% [5] coauthors
  
    
   
Sampling-based motion planning methods, while effective in high-dimensional spaces, often suffer from inefficiencies due to irregular sampling distributions, leading to suboptimal exploration of the configuration space. In this talk, we present an approach that enhances the efficiency of these methods by utilizing low-discrepancy distributions generated through Message-Passing Monte Carlo (MPMC). MPMC leverages Graph Neural Networks (GNNs) to generate point sets that uniformly cover the space, with uniformity assessed using the $\mathcal{L}_p$-discrepancy measure, which quantifies the irregularity of sample distributions. By improving the uniformity of the point sets, our approach significantly reduces computational overhead and the number of samples required for solving motion planning problems. Experimental results demonstrate that our method outperforms traditional sampling techniques in terms of planning efficiency.

\medskip

% If you would like to include references, please do so by creating a simple list numbered by [1], [2], [3], \ldots. See example below.
% Please do not use the \texttt{bibliography} environment or \texttt{bibtex} files.
% APA reference style is recommended.
% \begin{enumerate}
%  \item[{[1]}] Niederreiter, Harald (1992). {\it Random number generation and quasi-Monte Carlo methods}. Society for Industrial and Applied Mathematics (SIAM).
%  \item[{[2]}] L’Ecuyer, Pierre, \& Christiane Lemieux. (2002). Recent advances in randomized quasi-Monte Carlo methods. Modeling uncertainty: An examination of stochastic theory, methods, and applications, 419-474.
% \end{enumerate}

% Equations may be used if they are referenced. Please note that the equation numbers may be different (but will be cross-referenced correctly) in the final program book.
\end{talk}

\begin{talk}
  {An Empirical Evaluation of Robust Estimators for RQMC}%
% {Cost-Free improvements to the median estimator for Random-parameter QMC}% [1] talk title
  {Gregory Seljak}% [2] speaker name
  {Universit\'e de Montr\'eal}% [3] affiliations
  {gregory.de.salaberry.seljak@umontreal.ca}% [4] email
  {Pierre L'Ecuyer, Christiane Lemieux}% [5] coauthors
  {Computational Methods for Low-discrepancy Sampling and Applications}% [6] special session. 

\medskip

Randomized quasi-Monte Carlo (RQMC) traditionally takes a low-discrepancy (QMC) point set,
makes $r$ independent randomizations of it to obtain $r$ replicates of an unbiased RQMC estimator, 
then computes the average and variance of these $r$ estimates to obtain a final estimate 
and perhaps a confidence interval [4]. 
Some methods construct the points by optimizing refined figures-of-merit adapted to the considered integrand
and apply minimal randomization such as a random (digital) shift.  Other methods randomize
the parameters of the QMC point sets more extensively (e.g., the generating vectors or matrices).
The second kind of method is easier to apply because it requires much less knowledge of the integrand,
but bad parameter values may be drawn once in a while, leading to (rare) RQMC replicates having 
a large conditional variance that produce bad outliers.
To reduce the impact of such outliers, one approach studied recently is to use the median 
of the $r$ replicates instead of the mean as a final estimator [2, 3, 5, 6].
Other types of robust estimators could also be used in place of the median [1].
In this talk, we report extensive experiments that compare the mean square errors 
and convergence of various estimators (the mean, the median, and other robust estimators) 
defined in terms of $r$ RQMC replicates.
We also discuss the computation of confidence intervals for the mean when using such estimators.

{\list{[\arabic{enumi}]}{\settowidth\labelwidth{[5]}\leftmargin\labelwidth
  \advance\leftmargin\labelsep\usecounter{enumi}}
%
\item
E.~Gobet, M.~Lerasle, and D.~M{\'e}tivier.
Accelerated convergence of error quantiles using robust randomized
  quasi {Monte Carlo} methods.
\url{https://hal.science/hal-03631879}, 2024.

\item
T.~Goda and P.~L'Ecuyer.
Construction-free median quasi-{Monte Carlo} rules for function
  spaces with unspecified smoothness and general weights.
{\em {SIAM} Journal on Scientific Computing}, 44(4):A2765--A2788, 2022.

\item
T.~Goda, K.~Suzuki, and M.~Matsumoto.
A universal median quasi-{Monte Carlo} integration.
{\em {SIAM} Journal on Numerical Analysis}, 62(1):533--566, 2024.

\item
P.~L'Ecuyer, M.~Nakayama, A.~B. Owen, and B.~Tuffin.
Confidence intervals for randomized quasi-{Monte Carlo} estimators.
In {\em Proceedings of the 2023 Winter Simulation Conference}, pages
  445--456. IEEE Press, 2023.

\item
Z.~Pan and A.~B. Owen.
Super-polynomial accuracy of one dimensional randomized nets using
  the median of means.
{\em Mathematics of Computation}, 92(340):805--837, 2023.

\item
Z.~Pan and A.~B. Owen.
Super-polynomial accuracy of multidimensional randomized nets using
  the median of means.
{\em Mathematics of Computation}, 93(349):2265--2289, 2024.

}
\end{talk}
